\begin{thebibliography}{999}
\bibitem{Bekenstein1981} Bekenstein, J. D. Universal upper bound on the entropy-to-energy ratio for bounded systems. Phys. Rev. D 23, 287 (1981)
\bibitem{Bekenstein1973} Bekenstein, J. D. Black holes and entropy. Phys. Rev. D 7, 2333 (1973)
\bibitem{Jacobson1995} Jacobson, T. Thermodynamics of spacetime: the Einstein equation of state. Phys. Rev. Lett. 75, 1260 (1995)
\bibitem{RyuTakayanagi2006} Ryu, S. and Takayanagi, T. Holographic derivation of entanglement entropy from the anti-de Sitter space/conformal field theory correspondence. Phys. Rev. Lett. 96, 181602 (2006)
\bibitem{Jacobson2016} Jacobson, T. Entanglement equilibrium and the Einstein equation. Phys. Rev. Lett. 116, 201101 (2016)
\bibitem{Wheeler1990} Wheeler, J. A. Information, physics, quantum: the search for links. In Complexity, Entropy and the Physics of Information, ed. W. H. Zurek. Addison-Wesley (1990)
\bibitem{Hardy2001} Hardy, L. Quantum theory from five reasonable axioms. arXiv:quant-ph/0101012 (2001)
\bibitem{Chiribella2011} Chiribella, G., D'Ariano, G. M., and Perinotti, P. Informational derivation of quantum theory. Phys. Rev. A 84, 012311 (2011)
\bibitem{GeorgiGlashow1974} Georgi, H. and Glashow, S. L. Unity of all elementary-particle forces. Phys. Rev. Lett. 32, 438 (1974)
\bibitem{Horodecki1998} Horodecki, M., Horodecki, P., and Horodecki, R. Mixed-state entanglement and distillation: Is there a ``bound'' entanglement in nature? Phys. Rev. Lett. 80, 5239 (1998)
\bibitem{Horodecki2009} Horodecki, R., Horodecki, P., Horodecki, M., and Horodecki, K. Quantum entanglement. Rev. Mod. Phys. 81, 865 (2009)
\bibitem{Terhal2000} Terhal, B. M. Bell inequalities and the separability criterion. Phys. Lett. A 271, 319 (2000)
\bibitem{Peres1996} Peres, A. Separability criterion for density matrices. Phys. Rev. Lett. 77, 1413 (1996)
\bibitem{DiVincenzo2002} DiVincenzo, D. P., Leung, D. W., and Terhal, B. M. Quantum data hiding. IEEE Trans. Inf. Theory 48, 580 (2002)
\bibitem{Weinberg1987} Weinberg, S. Anthropic bound on the cosmological constant. Phys. Rev. Lett. 59, 2607 (1987)
\bibitem{Vafa2005} Vafa, C. The string landscape and the swampland. arXiv:hep-th/0509212 (2005)
\bibitem{BilsonThompson2005} Bilson-Thompson, S. O. A topological model of composite preons. arXiv:hep-ph/0503213 (2005)
\bibitem{Lisi2007} Lisi, A. G. An exceptionally simple theory of everything. arXiv:0711.0770 (2007)
\bibitem{Furey2016} Furey, C. Standard model physics from an algebra? PhD thesis, University of Waterloo, arXiv:1611.09182 (2015)
\bibitem{ChamseddineConnes1997} Chamseddine, A. H. and Connes, A. The spectral action principle. Commun. Math. Phys. 186, 731 (1997)
\bibitem{Eddington1946} Eddington, A. S. Fundamental Theory. Cambridge University Press (1946)
\bibitem{Zukowski1993} Zukowski, M., Zeilinger, A., Horne, M. A., and Ekert, A. K. ``Event-ready-detectors'' Bell experiment via entanglement swapping. Phys. Rev. Lett. 71, 4287 (1993)
\bibitem{Weyl1927} Weyl, H. Quantenmechanik und Gruppentheorie. Z. Phys. 46, 1 (1927)
\bibitem{Gottesman1997} Gottesman, D. Stabilizer codes and quantum error correction. PhD thesis, California Institute of Technology, arXiv:quant-ph/9705052 (1997)
\bibitem{WoottersZurek1982} Wootters, W. K. and Zurek, W. H. A single quantum cannot be cloned. Nature 299, 802 (1982)
\bibitem{Schrodinger1935} Schroedinger, E. Discussion of probability relations between separated systems. Proc. Cambridge Philos. Soc. 31, 555 (1935)
\bibitem{EPR1935} Einstein, A., Podolsky, B., and Rosen, N. Can quantum-mechanical description of physical reality be considered complete? Phys. Rev. 47, 777 (1935)
\bibitem{Hilbert1899} Hilbert, D. Grundlagen der Geometrie. Teubner (1899)
\bibitem{Mackey1963} Mackey, G. W. The Mathematical Foundations of Quantum Mechanics. Benjamin (1963)
\bibitem{Einstein1905} Einstein, A. Zur Elektrodynamik bewegter Koerper. Ann. Phys. 17, 891 (1905)
\bibitem{Wigner1960} Wigner, E. P. The unreasonable effectiveness of mathematics in the natural sciences. Commun. Pure Appl. Math. 13, 1 (1960)
\bibitem{Schmidt1907} Schmidt, E. Zur Theorie der linearen und nichtlinearen Integralgleichungen. Math. Ann. 63, 433 (1907)
\bibitem{Nielsen2000} Nielsen, M. A. and Chuang, I. L. Quantum Computation and Quantum Information. Cambridge University Press (2000)
\bibitem{Dirac1927} Dirac, P. A. M. The quantum theory of the emission and absorption of radiation. Proc. R. Soc. A 114, 243 (1927)
\bibitem{vonNeumann1932} von Neumann, J. Mathematische Grundlagen der Quantenmechanik. Springer, Berlin (1932)
\bibitem{Nielsen1999} Nielsen, M. A. Conditions for a class of entanglement transformations. Phys. Rev. Lett. 83, 436 (1999)
\bibitem{Ward1950} Ward, J. C. An identity in quantum electrodynamics. Phys. Rev. 78, 182 (1950)
\bibitem{BjorkenDrell1965} Bjorken, J. D. and Drell, S. D. Relativistic Quantum Fields. McGraw-Hill (1965)
\bibitem{Low1958} Low, F. E. Bremsstrahlung of very low-energy quanta in elementary particle collisions. Phys. Rev. 110, 974 (1958)
\bibitem{Weinberg1965} Weinberg, S. Infrared photons and gravitons. Phys. Rev. 140, B516 (1965)
\bibitem{Berry1984} Berry, M. V. Quantal phase factors accompanying adiabatic changes. Proc. R. Soc. A 392, 45 (1984)
\bibitem{Pancharatnam1956} Pancharatnam, S. Generalized theory of interference, and its applications: Part I. Coherent pencils. Proc. Indian Acad. Sci. A 44, 247 (1956)
\bibitem{Glashow1961} Glashow, S. L. Partial-symmetries of weak interactions. Nucl. Phys. 22, 579 (1961)
\bibitem{Weinberg1967} Weinberg, S. A Model of Leptons. Phys. Rev. Lett. 19, 1264 (1967)
\bibitem{Feynman1963} Feynman, R. P. Quantum theory of gravitation. Acta Phys. Polon. 24, 697 (1963)
\bibitem{BlochNordsieck1937} Bloch, F. and Nordsieck, A. Note on the radiation field of the electron. Phys. Rev. 52, 54 (1937)
\bibitem{Dyson1949} Dyson, F. J. The S matrix in quantum electrodynamics. Phys. Rev. 75, 1736 (1949)
\bibitem{Aharonov1959} Aharonov, Y. and Bohm, D. Significance of electromagnetic potentials in the quantum theory. Phys. Rev. 115, 485 (1959)
\bibitem{WuYang1975} Wu, T. T. and Yang, C. N. Concept of nonintegrable phase factors and global formulation of gauge fields. Phys. Rev. D 12, 3845 (1975)
\bibitem{Einstein1916} Einstein, A. Die Grundlage der allgemeinen Relativitaetstheorie. Ann. Phys. 49, 769 (1916)
\bibitem{Born1926} Born, M. Zur Quantenmechanik der Stossvorgaenge. Z. Phys. 37, 863 (1926)
\bibitem{Zurek2003} Zurek, W. H. Decoherence, einselection, and the quantum origins of the classical. Rev. Mod. Phys. 75, 715 (2003)
\bibitem{Misra1977} Misra, B. and Sudarshan, E. C. G. The Zeno's paradox in quantum theory. J. Math. Phys. 18, 756 (1977)
\bibitem{Schur1905} Schur, I. Neue Begruendung der Theorie der Gruppencharaktere. Sitzungsber. Preuss. Akad. Wiss. Berlin, 406 (1905)
\bibitem{Weyl1931} Weyl, H. The Theory of Groups and Quantum Mechanics. Methuen, London (1931)
\bibitem{Werner1989} Werner, R. F. Quantum states with Einstein-Podolsky-Rosen correlations admitting a hidden-variable model. Phys. Rev. A 40, 4277 (1989)
\bibitem{Jarlskog1985} Jarlskog, C. A basis independent formulation of the connection between quark mass matrices, CP violation and experiment. Z. Phys. C 29, 491 (1985)
\bibitem{Kobayashi1973} Kobayashi, M. and Maskawa, T. CP-Violation in the Renormalizable Theory of Weak Interaction. Prog. Theor. Phys. 49, 652 (1973)
\bibitem{Bohr1920} Bohr, N. Ueber die Serienspektra der Elemente. Z. Phys. 2, 423 (1920)
\bibitem{Ehrenfest1927} Ehrenfest, P. Bemerkung ueber die angenaeherte Gueltigkeit der klassischen Mechanik innerhalb der Quantenmechanik. Z. Phys. 45, 455 (1927)
\bibitem{AppelquistCarazzone1975} Appelquist, T. and Carazzone, J. Infrared singularities and massive fields. Phys. Rev. D 11, 2856 (1975)
\bibitem{Bell1964} Bell, J. S. On the Einstein Podolsky Rosen paradox. Physics 1, 195 (1964)
\bibitem{Tsirelson1980} Cirel'son, B. S. Quantum generalizations of Bell's inequality. Lett. Math. Phys. 4, 93 (1980)
\bibitem{CGLMP2002} Collins, D., Gisin, N., Linden, N., Massar, S., and Popescu, S. Bell inequalities for arbitrarily high-dimensional systems. Phys. Rev. Lett. 88, 040404 (2002)
\bibitem{vonNeumann1931} von Neumann, J. Die Eindeutigkeit der Schroedingerschen Operatoren. Math. Ann. 104, 570 (1931)
\bibitem{Weyl1928} Weyl, H. Gruppentheorie und Quantenmechanik. Hirzel (1928)
\bibitem{Dirac1930} Dirac, P. A. M. The Principles of Quantum Mechanics. Clarendon Press, Oxford (1930)
\bibitem{Hein2004} Hein, M., Eisert, J., and Briegel, H. J. Multiparty entanglement in graph states. Phys. Rev. A 69, 062311 (2004)
\bibitem{Yukawa1935} Yukawa, H. On the interaction of elementary particles. I. Proc. Phys.-Math. Soc. Japan 17, 48 (1935)
\bibitem{Gleason1957} Gleason, A. M. Measures on the closed subspaces of a Hilbert space. J. Math. Mech. 6, 885 (1957)
\bibitem{Zurek2003envariance} Zurek, W. H. Environment-assisted invariance, entanglement, and probabilities in quantum physics. Phys. Rev. Lett. 90, 120404 (2003)
\bibitem{Zurek2005} Zurek, W. H. Probabilities from entanglement, Born's rule $p_k=|\psi_k|^2$ from envariance. Phys. Rev. A 71, 052105 (2005)
\bibitem{Robertson1929} Robertson, H. P. The uncertainty principle. Phys. Rev. 34, 163 (1929)
\bibitem{Peres1993} Peres, A. Quantum Theory: Concepts and Methods. Kluwer (1993)
\bibitem{Raussendorf2001} Raussendorf, R. and Briegel, H. J. A one-way quantum computer. Phys. Rev. Lett. 86, 5188 (2001)
\bibitem{Pauli1940} Pauli, W. The connection between spin and statistics. Phys. Rev. 58, 716 (1940)
\bibitem{Stone1932} Stone, M. H. On one-parameter unitary groups in Hilbert space. Ann. Math. 33, 643 (1932)
\bibitem{Dirac1958} Dirac, P. A. M. The Principles of Quantum Mechanics, fourth edition. Oxford University Press (1958)
\bibitem{Bass1992} Bass, H. The Ihara-Selberg zeta function of a tree lattice. Int. J. Math. 3, 717 (1992)
\bibitem{Feynman1948} Feynman, R. P. Space-Time Approach to Non-Relativistic Quantum Mechanics. Rev. Mod. Phys. 20, 367 (1948)
\bibitem{Peskin1995} Peskin, M. E. and Schroeder, D. V. An Introduction to Quantum Field Theory. Addison-Wesley (1995)
\bibitem{Wigner1932} Wigner, E. P. On the quantum correction for thermodynamic equilibrium. Phys. Rev. 40, 749 (1932)
\bibitem{Sakurai1994} Sakurai, J. J. Modern Quantum Mechanics, revised edition. Addison-Wesley (1994)
\bibitem{Holevo1973} Holevo, A. S. Bounds for the quantity of information transmitted by a quantum communication channel. Probl. Inf. Transm. 9, 177 (1973)
\bibitem{Minkowski1908} Minkowski, H. Raum und Zeit. Phys. Z. 10, 104 (1909)
\bibitem{Wilson1974} Wilson, K. G. Confinement of quarks. Phys. Rev. D 10, 2445 (1974)
\bibitem{Feynman1982} Feynman, R. P. Simulating physics with computers. Int. J. Theor. Phys. 21, 467 (1982)
\bibitem{tHooft2016} 't Hooft, G. The Cellular Automaton Interpretation of Quantum Mechanics (Springer, Cham, 2016)
\bibitem{LSZ1955} Lehmann, H., Symanzik, K., and Zimmermann, W. Zur Formulierung quantisierter Feldtheorien. Nuovo Cim. 1, 205 (1955)
\bibitem{Hamermesh1962} Hamermesh, M. Group Theory and Its Application to Physical Problems. Addison-Wesley (1962)
\bibitem{Wigner1959} Wigner, E. P. Group Theory and Its Application to the Quantum Mechanics of Atomic Spectra. Academic Press (1959)
\bibitem{Weinberg1995} Weinberg, S. The Quantum Theory of Fields, Volume I: Foundations. Cambridge University Press (1995)
\bibitem{HardyWright1979} Hardy, G. H. and Wright, E. M. An Introduction to the Theory of Numbers, fifth edition. Oxford University Press (1979)
\bibitem{Wigner1939} Wigner, E. P. On unitary representations of the inhomogeneous Lorentz group. Ann. Math. 40, 149 (1939)
\bibitem{Dirac1928} Dirac, P. A. M. The quantum theory of the electron. Proc. R. Soc. Lond. A 117, 610 (1928)
\bibitem{AtiyahBottShapiro1964} Atiyah, M. F., Bott, R., and Shapiro, A. Clifford modules. Topology 3, Suppl. 1, 3 (1964)
\bibitem{Lee1956} Lee, T. D. and Yang, C. N. Question of Parity Conservation in Weak Interactions. Phys. Rev. 104, 254 (1956)
\bibitem{Luders1957} Lueders, G. Proof of the TCP theorem. Ann. Phys. (N.Y.) 2, 1 (1957)
\bibitem{Jost1957} Jost, R. Eine Bemerkung zum CTP Theorem. Helv. Phys. Acta 30, 409 (1957)
\bibitem{StreaterWightman1964} Streater, R. F. and Wightman, A. S. PCT, Spin and Statistics, and All That. Benjamin (1964)
\bibitem{Wu1957} Wu, C. S., Ambler, E., Hayward, R. W., Hoppes, D. D., and Hudson, R. P. Experimental Test of Parity Conservation in Beta Decay. Phys. Rev. 105, 1413 (1957)
\bibitem{Luders1954} Lueders, G. On the equivalence of invariance under time reversal and under particle-antiparticle conjugation for relativistic field theories. Dan. Mat. Fys. Medd. 28, 5 (1954)
\bibitem{Pauli1955} Pauli, W. Exclusion Principle, Lorentz Group and Reflection of Space-Time and Charge. In Niels Bohr and the Development of Physics, p. 30 (Pergamon, London, 1955)
\bibitem{Landau1937} Landau, L. D. On the theory of phase transitions. Zh. Eksp. Teor. Fiz. 7, 19 (1937)
\bibitem{Noether1918} Noether, E. Invariante Variationsprobleme. Nachr. Ges. Wiss. Goettingen, Math.-Phys. Kl. 1918, 235 (1918)
\bibitem{Kemeny1960} Kemeny, J. G. and Snell, J. L. Finite Markov Chains. Van Nostrand (1960)
\bibitem{YangMills1954} Yang, C. N. and Mills, R. L. Conservation of isotopic spin and isotopic gauge invariance. Phys. Rev. 96, 191 (1954)
\bibitem{Baez2010} Baez, J. C. and Huerta, J. The algebra of grand unified theories. Bull. Amer. Math. Soc. 47, 483 (2010)
\bibitem{Hall2015} Hall, B. C. Lie groups, Lie algebras, and representations. Springer (2015)
\bibitem{Humphreys1972} Humphreys, J. E. Introduction to Lie algebras and representation theory. Springer (1972)
\bibitem{FritzschGellMannLeutwyler1973} Fritzsch, H., Gell-Mann, M., and Leutwyler, H. Advantages of the color octet gluon picture. Phys. Lett. B 47, 365 (1973)
\bibitem{Bouchiat1972} Bouchiat, C., Iliopoulos, J., and Meyer, P. An anomaly-free version of Weinberg's model. Phys. Lett. B 38, 519 (1972)
\bibitem{GellMann1964} Gell-Mann, M. A schematic model of baryons and mesons. Phys. Lett. 8, 214 (1964)
\bibitem{Feynman1949} Feynman, R. P. Space-Time Approach to Quantum Electrodynamics. Phys. Rev. 76, 769 (1949)
\bibitem{Schwinger1951} Schwinger, J. On the Green's functions of quantized fields. I. Proc. Natl. Acad. Sci. 37, 452 (1951)
\bibitem{Brillouin1946} Brillouin, L. Wave Propagation in Periodic Structures (McGraw-Hill, New York, 1946)
\bibitem{GellMannNishijima1953} Gell-Mann, M. Isotopic spin and new unstable particles. Phys. Rev. 92, 833 (1953); Nakano, T. and Nishijima, K. Charge independence for V-particles. Prog. Theor. Phys. 10, 581 (1953)
\bibitem{Feinberg1987} Feinberg, M. Chemical reaction network structure and the stability of complex isothermal reactors, I. The deficiency zero and deficiency one theorems. Chem. Eng. Sci. 42, 2229 (1987)
\bibitem{Takahashi1957} Takahashi, Y. On the generalized Ward identity. Nuovo Cimento 6, 371 (1957)
\bibitem{Schwinger1949} Schwinger, J. Quantum electrodynamics. III. The electromagnetic properties of the electron: radiative corrections to scattering. Phys. Rev. 76, 790 (1949)
\bibitem{Adler1969} Adler, S. L. Axial-vector vertex in spinor electrodynamics. Phys. Rev. 177, 2426 (1969)
\bibitem{BellJackiw1969} Bell, J. S. and Jackiw, R. A PCAC puzzle: $\pi^0\to\gamma\gamma$ in the $\sigma$-model. Nuovo Cimento A 60, 47 (1969)
\bibitem{Schwinger1948} Schwinger, J. On Quantum-Electrodynamics and the Magnetic Moment of the Electron. Phys. Rev. 73, 416 (1948)
\bibitem{Kinoshita1990} Kinoshita, T. (ed.). Quantum Electrodynamics. World Scientific, Singapore (1990)
\bibitem{Petermann1957} Petermann, A. Fourth order magnetic moment of the electron. Helv. Phys. Acta 30, 407 (1957)
\bibitem{Higgs1964} Higgs, P. W. Broken symmetries and the masses of gauge bosons. Phys. Rev. Lett. 13, 508 (1964)
\bibitem{Maxwell1865} Maxwell, J. C. A dynamical theory of the electromagnetic field. Phil. Trans. R. Soc. Lond. 155, 459 (1865)
\bibitem{Einstein1917} Einstein, A. Kosmologische Betrachtungen zur allgemeinen Relativitaetstheorie. Sitzungsber. Preuss. Akad. Wiss. 1917, 142 (1917)
\bibitem{Einstein1907} Einstein, A. Ueber das Relativitaetsprinzip und die aus demselben gezogenen Folgerungen. Jahrbuch der Radioaktivitaet und Elektronik 4, 411 (1907)
\bibitem{Will2014} Will, C. M. The Confrontation between General Relativity and Experiment. Living Rev. Relativ. 17, 4 (2014)
\bibitem{Wald1984} Wald, R. M. General Relativity. University of Chicago Press (1984)
\bibitem{Hawking1975} Hawking, S. W. Particle Creation by Black Holes. Commun. Math. Phys. 43, 199 (1975)
\bibitem{Bousso2002} Bousso, R. The Holographic Principle. Rev. Mod. Phys. 74, 825 (2002)
\bibitem{Hubble1929} Hubble, E. A relation between distance and radial velocity among extra-galactic nebulae. Proc. Natl. Acad. Sci. 15, 168 (1929)
\bibitem{Schwarzschild1916} Schwarzschild, K. Ueber das Gravitationsfeld eines Massenpunktes nach der Einsteinschen Theorie. Sitzungsber. Preuss. Akad. Wiss. Berlin, 189 (1916)
\bibitem{Alpher1948} Alpher, R. A., Bethe, H., and Gamow, G. The Origin of Chemical Elements. Phys. Rev. 73, 803 (1948)
\bibitem{Friedmann1922} Friedmann, A. Ueber die Kruemmung des Raumes. Z. Phys. 10, 377 (1922)
\bibitem{HawkingEllis1973} Hawking, S. W. and Ellis, G. F. R. The Large Scale Structure of Space-Time. Cambridge University Press (1973)
\bibitem{Newton1687} Newton, I. Philosophiae Naturalis Principia Mathematica. London (1687)
\bibitem{Chung1997} Chung, F. R. K. Spectral Graph Theory. American Mathematical Society (1997)
\bibitem{Gibbs1902} Gibbs, J. W. Elementary Principles in Statistical Mechanics. Scribner (1902)
\bibitem{Planck1899} Planck, M. Ueber irreversible Strahlungsvorgaenge. Sitzungsber. Preuss. Akad. Wiss. 5, 440 (1899)
\bibitem{Wilczek2005} Wilczek, F. On Absolute Units, I: Choices. Phys. Today 58, 12 (2005)
\bibitem{Sommerfeld1916} Sommerfeld, A. Zur Quantentheorie der Spektrallinien. Ann. Phys. 356, 1 (1916)
\bibitem{Mohr2025} Mohr, P. J. et al. CODATA recommended values of the fundamental physical constants: 2022. Rev. Mod. Phys. 97, 025002 (2025)
\bibitem{Weyl1929} Weyl, H. Elektron und Gravitation. Z. Phys. 56, 330 (1929)
\bibitem{Gasser1982} Gasser, J. and Leutwyler, H. Quark masses. Phys. Rep. 87, 77 (1982)
\bibitem{Hanneke2008} Hanneke, D., Fogwell, S., and Gabrielse, G. New Measurement of the Electron Magnetic Moment and the Fine Structure Constant. Phys. Rev. Lett. 100, 120801 (2008)
\bibitem{Horodecki1999} Horodecki, M. and Horodecki, P. Reduction criterion of separability and limits for a class of distillation protocols. Phys. Rev. A 59, 4206 (1999)
\bibitem{Tiesinga2021} Tiesinga, E., Mohr, P. J., Newell, D. B., and Taylor, B. N. CODATA recommended values of the fundamental physical constants: 2018. Rev. Mod. Phys. 93, 025010 (2021)
\bibitem{Aoyama2018} Aoyama, T., Kinoshita, T., and Nio, M. Revised and Improved Value of the QED Tenth-Order Electron Anomalous Magnetic Moment. Phys. Rev. D 97, 036001 (2018)
\bibitem{Bohr1913} Bohr, N. On the constitution of atoms and molecules. Philos. Mag. 26, 1 (1913)
\bibitem{Lamb1947} Lamb, W. E. and Retherford, R. C. Fine Structure of the Hydrogen Atom by a Microwave Method. Phys. Rev. 72, 241 (1947)
\bibitem{Ramsey1990} Ramsey, N. F. Experiments with separated oscillatory fields and hydrogen masers. Rev. Mod. Phys. 62, 541 (1990)
\bibitem{Fermi1930} Fermi, E. Ueber die magnetischen Momente der Atomkerne. Z. Phys. 60, 320 (1930)
\bibitem{Ramsey1965} Kleppner, D., Berg, H. C., Crampton, S. B., Ramsey, N. F., Vessot, R. F. C., Peters, H. E., and Vanier, J. Hydrogen-Maser Principles and Techniques. Phys. Rev. 138, A972 (1965)
\bibitem{Casimir1948} Casimir, H. B. G. On the attraction between two perfectly conducting plates. Proc. K. Ned. Akad. Wet. 51, 793 (1948)
\bibitem{Bethe1947} Bethe, H. A. The Electromagnetic Shift of Energy Levels. Phys. Rev. 72, 339 (1947)
\bibitem{Dyson1952} Dyson, F. J. Divergence of Perturbation Theory in Quantum Electrodynamics. Phys. Rev. 85, 631 (1952)
\bibitem{GellMann1954} Gell-Mann, M. and Low, F. E. Quantum Electrodynamics at Small Distances. Phys. Rev. 95, 1300 (1954)
\bibitem{Serber1935} Serber, R. Linear Modifications in the Maxwell Field Equations. Phys. Rev. 48, 49 (1935)
\bibitem{Gross1973} Gross, D. J. and Wilczek, F. Ultraviolet behavior of non-abelian gauge theories. Phys. Rev. Lett. 30, 1343 (1973)
\bibitem{Politzer1973} Politzer, H. D. Reliable Perturbative Results for Strong Interactions? Phys. Rev. Lett. 30, 1346 (1973)
\bibitem{BetheHeitler1934} Bethe, H. and Heitler, W. On the stopping of fast particles and on the creation of positive electrons. Proc. R. Soc. Lond. A 146, 83 (1934)
\bibitem{Pontecorvo1958} Pontecorvo, B. Inverse beta processes and nonconservation of lepton charge. Sov. Phys. JETP 7, 172 (1958)
\bibitem{Maki1962} Maki, Z., Nakagawa, M., and Sakata, S. Remarks on the unified model of elementary particles. Prog. Theor. Phys. 28, 870 (1962)
\bibitem{Koide1983} Koide, Y. A fermion-boson composite model of quarks and leptons. Phys. Lett. B 120, 161 (1983)
\bibitem{Jackson1998} Jackson, J. D. Classical Electrodynamics, 3rd ed. Wiley (1998)
\bibitem{Salam1968} Salam, A. Weak and Electromagnetic Interactions. In Elementary Particle Theory, ed. N. Svartholm (Almqvist and Wiksell, 1968), p. 367
\bibitem{Cabibbo1963} Cabibbo, N. Unitary symmetry and leptonic decays. Phys. Rev. Lett. 10, 531 (1963)
\bibitem{Chew1962} Chew, G. F. and Frautschi, S. C. Regge trajectories and the principle of maximum strength for strong interactions. Phys. Rev. Lett. 8, 41 (1962)
\bibitem{Mathieu2009} Mathieu, V., Kochelev, N., and Vento, V. The physics of glueballs. Int. J. Mod. Phys. E 18, 1 (2009)
\bibitem{Weinberg1977} Weinberg, S. The problem of mass. Trans. N. Y. Acad. Sci. 38, 185 (1977)
\bibitem{Perl1975} Perl, M. L. et al. Evidence for Anomalous Lepton Production in e+e- Annihilation. Phys. Rev. Lett. 35, 1489 (1975)
\bibitem{Avigad2021} Avigad, J. and Massot, P. Mathematics in Lean. Lean community (2021)
\bibitem{GellMann1962} Gell-Mann, M. Symmetries of baryons and mesons. Phys. Rev. 125, 1067 (1962)
\bibitem{Okubo1962} Okubo, S. Note on unitary symmetry in strong interactions. Prog. Theor. Phys. 27, 949 (1962)
\bibitem{BreitWigner1936} Breit, G. and Wigner, E. Capture of slow neutrons. Phys. Rev. 49, 519 (1936)
\bibitem{Fermi1934} Fermi, E. Versuch einer Theorie der beta-Strahlen. Z. Phys. 88, 161 (1934)
\bibitem{WeisskopfWigner1930} Weisskopf, V. and Wigner, E. Berechnung der natuerlichen Linienbreite auf Grund der Diracschen Lichttheorie. Z. Phys. 63, 54 (1930)
\bibitem{Weizsaecker1935} von Weizsaecker, C. F. Zur Theorie der Kernmassen. Z. Phys. 96, 431 (1935)
\bibitem{BetheBacher1936} Bethe, H. A. and Bacher, R. F. Nuclear physics A. Stationary states of nuclei. Rev. Mod. Phys. 8, 82 (1936)
\bibitem{Bethe1935} Bethe, H. A. and Peierls, R. Quantum theory of the diplon. Proc. R. Soc. Lond. A 148, 146 (1935)
\bibitem{Weinberg1991} Weinberg, S. Effective chiral Lagrangians for nucleon-pion interactions and nuclear forces. Nucl. Phys. B 363, 3 (1991)
\bibitem{Mayer1949} Mayer, M. G. On closed shells in nuclei II. Phys. Rev. 75, 1969 (1949)
\bibitem{Clausius1865} Clausius, R. Ueber verschiedene fuer die Anwendung bequeme Formen der Hauptgleichungen der mechanischen Waermetheorie. Annalen der Physik und Chemie 125, 353 (1865)
\bibitem{Boltzmann1872} Boltzmann, L. Weitere Studien ueber das Waermegleichgewicht unter Gasmolekuelen. Sitzungsberichte der Kaiserlichen Akademie der Wissenschaften Wien 66, 275 (1872)
\bibitem{Callen1985} Callen, H. B. Thermodynamics and an Introduction to Thermostatistics, second edition. Wiley, New York (1985)
\bibitem{Bennett1996} Bennett, C. H., DiVincenzo, D. P., Smolin, J. A., and Wootters, W. K. Mixed-state entanglement and quantum error correction. Phys. Rev. A 54, 3824 (1996)
\bibitem{Onsager1944} Onsager, L. Crystal statistics. I. A two-dimensional model with an order-disorder transition. Phys. Rev. 65, 117 (1944)
\bibitem{Jarzynski1997} Jarzynski, C. Nonequilibrium equality for free energy differences. Phys. Rev. Lett. 78, 2690 (1997)
\bibitem{Crooks1999} Crooks, G. E. Entropy production fluctuation theorem and the nonequilibrium work relation for free energy differences. Phys. Rev. E 60, 2721 (1999)
\bibitem{Onsager1931} Onsager, L. Reciprocal relations in irreversible processes. I. Phys. Rev. 37, 405 (1931)
\bibitem{AshcroftMermin1976} Ashcroft, N. W. and Mermin, N. D. Solid State Physics. Holt, Rinehart and Winston (1976)
\bibitem{Kittel2005} Kittel, C. Introduction to Solid State Physics, eighth edition. Wiley (2005)
\bibitem{Bloch1929} Bloch, F. Ueber die Quantenmechanik der Elektronen in Kristallgittern. Z. Phys. 52, 555 (1929)
\bibitem{KronigPenney1931} Kronig, R. de L. and Penney, W. G. Quantum mechanics of electrons in crystal lattices. Proc. R. Soc. A 130, 499 (1931)
\bibitem{Peierls1955} Peierls, R. E. Quantum Theory of Solids. Oxford University Press (1955)
\bibitem{SuSchriefferHeeger1979} Su, W. P., Schrieffer, J. R., and Heeger, A. J. Solitons in polyacetylene. Phys. Rev. Lett. 42, 1698 (1979)
\bibitem{Mendeleev1869} Mendeleev, D. Ueber die Beziehungen der Eigenschaften zu den Atomgewichten der Elemente. Z. Chem. 12, 405 (1869)
\bibitem{Pauli1925} Pauli, W. Ueber den Zusammenhang des Abschlusses der Elektronengruppen im Atom mit der Komplexstruktur der Spektren. Z. Phys. 31, 765 (1925)
\bibitem{Schrodinger1926} Schroedinger, E. Quantisierung als Eigenwertproblem. Ann. Phys. 384, 361 (1926)
\bibitem{Hartree1928} Hartree, D. R. The wave mechanics of an atom with a non-Coulomb central field. Part I. Theory and methods. Proc. Camb. Philos. Soc. 24, 89 (1928)
\bibitem{Fock1930} Fock, V. Naeherungsmethode zur Loesung des quantenmechanischen Mehrkoerperproblems. Z. Phys. 61, 126 (1930)
\bibitem{Gottesman1998} Gottesman, D. The Heisenberg Representation of Quantum Computers. In Group22: Proceedings of the XXII International Colloquium on Group Theoretical Methods in Physics. International Press, Cambridge, MA, 32 (1998)
\bibitem{BravyiKitaev2005} Bravyi, S. and Kitaev, A. Universal quantum computation with ideal Clifford gates and noisy ancillas. Phys. Rev. A 71, 022316 (2005)
\bibitem{Landauer1961} Landauer, R. Irreversibility and heat generation in the computing process. IBM J. Res. Dev. 5, 183 (1961)
\bibitem{BornFock1928} Born, M. and Fock, V. Beweis des Adiabatensatzes. Z. Phys. 51, 165 (1928)
\bibitem{Simon1983} Simon, B. Holonomy, the quantum adiabatic theorem, and Berry's phase. Phys. Rev. Lett. 51, 2167 (1983)
\bibitem{Weinberg1996} Weinberg, S. The Quantum Theory of Fields, Volume II: Modern Applications. Cambridge University Press (1996)
\bibitem{Navas2024} Navas, S. et al. (Particle Data Group). Review of particle physics. Phys. Rev. D 110, 030001 (2024)
\bibitem{Weinberg1989} Weinberg, S. The cosmological constant problem. Rev. Mod. Phys. 61, 1 (1989)
\bibitem{Nishijima1955} Nishijima, K. Charge independence theory of V particles. Prog. Theor. Phys. 13, 285 (1955)
\bibitem{Griffiths2008} Griffiths, D. Introduction to Elementary Particles. Wiley-VCH, Weinheim, second edition (2008)
\bibitem{Chadwick1932} Chadwick, J. Possible existence of a neutron. Nature 129, 312 (1932)
\bibitem{Heisenberg1932} Heisenberg, W. Ueber den Bau der Atomkerne. I. Z. Phys. 77, 1 (1932)
\bibitem{Zwicky1933} Zwicky, F. Die Rotverschiebung von extragalaktischen Nebeln. Helv. Phys. Acta 6, 110 (1933)
\bibitem{LeeWeinberg1977} Lee, B. W. and Weinberg, S. Cosmological lower bound on heavy-neutrino masses. Phys. Rev. Lett. 39, 165 (1977)
\bibitem{FeynmanGellMann1958} Feynman, R. P. and Gell-Mann, M. Theory of the Fermi interaction. Phys. Rev. 109, 193 (1958)
\bibitem{Sakharov1967} Sakharov, A. D. Violation of CP invariance, C asymmetry, and baryon asymmetry of the universe. JETP Lett. 5, 24 (1967)
\bibitem{Christenson1964} Christenson, J. H., Cronin, J. W., Fitch, V. L., and Turlay, R. Evidence for the $2\pi$ decay of the $K_2^0$ meson. Phys. Rev. Lett. 13, 138 (1964)
\bibitem{Einstein1918} Einstein, A. Ueber Gravitationswellen. Sitzungsber. Preuss. Akad. Wiss. Berlin 1918, 154 (1918)
\bibitem{vonNeumann1966} von Neumann, J. Theory of Self-Reproducing Automata. University of Illinois Press (1966)
\bibitem{Rozenberg1997} Rozenberg, G. (ed.) Handbook of Graph Grammars and Computing by Graph Transformation, vol. 1. World Scientific (1997)
\bibitem{Newman1942} Newman, M. H. A. On theories with a combinatorial definition of equivalence. Ann. Math. 43, 223 (1942)
\bibitem{LiebRobinson1972} Lieb, E. H. and Robinson, D. W. The finite group velocity of quantum spin systems. Commun. Math. Phys. 28, 251 (1972)
\bibitem{Turing1937} Turing, A. M. On computable numbers, with an application to the Entscheidungsproblem. Proc. London Math. Soc. 42, 230 (1937)
\bibitem{Wolfram2002} Wolfram, S. A New Kind of Science. Wolfram Media (2002)
\bibitem{Lindblad1976} Lindblad, G. On the generators of quantum dynamical semigroups. Commun. Math. Phys. 48, 119 (1976)
\bibitem{Zurek1981} Zurek, W. H. Pointer basis of quantum apparatus: into what mixture does the wave packet collapse? Phys. Rev. D 24, 1516 (1981)
\bibitem{OllivierPoulinZurek2004} Ollivier, H., Poulin, D., and Zurek, W. H. Objective properties from subjective quantum states: environment as a witness. Phys. Rev. Lett. 93, 220401 (2004)
\bibitem{Zurek2009} Zurek, W. H. Quantum Darwinism. Nat. Phys. 5, 181 (2009)
\bibitem{tHooftVeltman1972} 't Hooft, G. and Veltman, M. Regularization and renormalization of gauge fields. Nucl. Phys. B 44, 189 (1972)
\bibitem{Weinberg1973} Weinberg, S. New approach to the renormalization group. Phys. Rev. D 8, 3497 (1973)
\bibitem{GellMannNeeman1964} Gell-Mann, M. and Ne'eman, Y. The Eightfold Way. W. A. Benjamin, New York (1964)
\bibitem{Zweig1964} Zweig, G. An SU(3) model for strong interaction symmetry and its breaking. CERN Report 8182/TH.401 (1964)
\bibitem{Workman2022} Workman, R. L. and others (Particle Data Group). Review of particle physics. Prog. Theor. Exp. Phys. 2022, 083C01 (2022)
\bibitem{Cavendish1798} Cavendish, H. Experiments to Determine the Density of the Earth. Phil. Trans. R. Soc. Lond. 88, 469 (1798)
\bibitem{Aoyama2020} Aoyama, T. et al. The anomalous magnetic moment of the muon in the Standard Model. Phys. Rep. 887, 1 (2020)
\bibitem{Majorana1937} Majorana, E. Teoria simmetrica dell'elettrone e del positrone. Nuovo Cimento 14, 171 (1937)
\bibitem{Weinberg1979} Weinberg, S. Baryon- and lepton-nonconserving processes. Phys. Rev. Lett. 43, 1566 (1979)
\bibitem{Hasert1973} Hasert, F. J. et al. (Gargamelle Collaboration). Observation of neutrino-like interactions without muon or electron in the Gargamelle neutrino experiment. Phys. Lett. B 46, 138 (1973)
\bibitem{Michel1950} Michel, L. Interaction between four half-spin particles and the decay of the mu-meson. Proc. Phys. Soc. A 63, 514 (1950)
\bibitem{NambuJonaLasinio1961} Nambu, Y. and Jona-Lasinio, G. Dynamical model of elementary particles based on an analogy with superconductivity. I. Phys. Rev. 122, 345 (1961)
\bibitem{Goldstone1961} Goldstone, J. Field theories with superconductor solutions. Nuovo Cimento 19, 154 (1961)
\bibitem{tHooft1976} 't Hooft, G. Symmetry breaking through Bell-Jackiw anomalies. Phys. Rev. Lett. 37, 8 (1976)
\bibitem{PecceiQuinn1977} Peccei, R. D. and Quinn, H. R. CP conservation in the presence of pseudoparticles. Phys. Rev. Lett. 38, 1440 (1977)
\bibitem{Gamow1928} Gamow, G. Zur Quantentheorie des Atomkernes. Z. Phys. 51, 204 (1928)
\bibitem{Gurney1928} Gurney, R. W. and Condon, E. U. Wave Mechanics and Radioactive Disintegration. Nature 122, 439 (1928)
\bibitem{Geiger1911} Geiger, H. and Nuttall, J. M. The ranges of the alpha particles from various radioactive substances and a relation between range and period of transformation. Phil. Mag. 22, 613 (1911)
\bibitem{BohrWheeler1939} Bohr, N. and Wheeler, J. A. The mechanism of nuclear fission. Phys. Rev. 56, 426 (1939)
\bibitem{Bethe1939} Bethe, H. A. Energy production in stars. Phys. Rev. 55, 434 (1939)
\bibitem{Chandrasekhar1931} Chandrasekhar, S. The maximum mass of ideal white dwarfs. Astrophys. J. 74, 81 (1931)
\bibitem{Thouless1982} Thouless, D. J., Kohmoto, M., Nightingale, M. P., and den Nijs, M. Quantized Hall Conductance in a Two-Dimensional Periodic Potential. Phys. Rev. Lett. 49, 405 (1982)
\bibitem{Klitzing1980} von Klitzing, K., Dorda, G., and Pepper, M. New method for high-accuracy determination of the fine-structure constant based on quantized Hall resistance. Phys. Rev. Lett. 45, 494 (1980)
\bibitem{deMoura2021} de Moura, L. and Ullrich, S. The Lean 4 theorem prover and programming language. In Automated Deduction - CADE 28. Lecture Notes in Computer Science 12699, 625 (2021)
\end{thebibliography}
